\begin{titlepage}
\centering
{\bfseries FACUMINAS\par}
{\bfseries Pós-graduação Lato Sensu em Educação Especial e Inclusiva, Neuropsicopedagogia Institucional e Clínica\par}
\vspace{3.5cm}
{\bfseries MARIA DO SOCORRO ARAUJO MARTINS LEITE\par}
\vspace{3.5cm}
{\bfseries A INCLUSÃO ESCOLAR DE ESTUDANTES COM TRANSTORNO DO ESPECTRO AUTISTA NOS ANOS INICIAIS: POLÍTICAS, PRÁTICAS PEDAGÓGICAS E FORMAÇÃO DOCENTE\par}
\vfill
{\bfseries Brasília\par}
{\bfseries 2026\par}
\end{titlepage}

\begin{titlepage}
\centering
{\bfseries MARIA DO SOCORRO ARAUJO MARTINS LEITE\par}
\vspace{3cm}
{\bfseries A INCLUSÃO ESCOLAR DE ESTUDANTES COM TRANSTORNO DO ESPECTRO AUTISTA NOS ANOS INICIAIS: POLÍTICAS, PRÁTICAS PEDAGÓGICAS E FORMAÇÃO DOCENTE\par}
\vspace{2.5cm}
\begin{flushright}
\begin{minipage}{9cm}
\small
Trabalho de Conclusão de Curso apresentado à Pós-graduação Lato Sensu em Educação Especial e Inclusiva, \mbox{Neuropsicopedagogia} Institucional e Clínica, da \mbox{FACUMINAS}, como requisito parcial para obtenção do título de Especialista em Educação Especial e Inclusiva, \mbox{Neuropsicopedagogia} Institucional e Clínica.\par
\end{minipage}
\end{flushright}
\vfill
{\bfseries Brasília\par}
{\bfseries 2026\par}
\end{titlepage}

\begin{center}
\textbf{RESUMO}
\end{center}

\begingroup
\setstretch{1.0}
\noindent
Este trabalho analisa como políticas educacionais, práticas pedagógicas e formação docente se articulam na inclusão escolar de estudantes com Transtorno do Espectro Autista (TEA) nos anos iniciais do ensino fundamental. Trata-se de revisão narrativa da literatura nacional, cujo corpus reúne oito obras com DOI ativo, publicadas entre 2013 e 2022, e cinco documentos legais. A pergunta que orienta a investigação é: como esses três eixos se articulam, na literatura, para viabilizar a participação e a aprendizagem de estudantes com TEA nos anos iniciais? A análise foi organizada em três eixos — políticas de inclusão, concepções sobre o TEA e práticas pedagógicas — e em uma síntese integrativa. Os resultados indicam que o avanço normativo não foi acompanhado, na mesma medida, de condições escolares efetivas. Persistem formação inicial insuficiente, escassez de recursos e de apoio especializado, rigidez curricular e barreiras atitudinais. Como estratégias, a literatura destaca adaptação curricular, ensino estruturado, comunicação alternativa, tecnologias assistivas, mediação do brincar, trabalho colaborativo e participação da família. Conclui-se que a inclusão não se realiza pela matrícula, mas pela reorganização estrutural, pedagógica e atitudinal da escola. O estudo é limitado por ser revisão narrativa, sem busca sistemática, e por não avaliar a qualidade metodológica das obras. Aponta-se a necessidade de pesquisas empíricas sobre práticas de sala de aula, formação docente e trajetórias de estudantes com TEA.

\vspace{1em}
\noindent\textbf{Palavras-chave:} Educação inclusiva. Transtorno do Espectro Autista. Anos iniciais. Formação docente. Práticas pedagógicas.
\endgroup

\clearpage

\begin{center}
\textbf{ABSTRACT}
\end{center}

\begingroup
\setstretch{1.0}
\noindent
This study analyzes how educational policies, pedagogical practices and teacher education articulate in the school inclusion of students with Autism Spectrum Disorder (ASD) in the early years of elementary school. It is a narrative review of Brazilian literature whose corpus comprises eight works with active DOI, published between 2013 and 2022, and five legal documents. The research question is: how do these three axes articulate, in the literature, to enable the participation and learning of students with ASD in the early years? The analysis was organized into three axes — inclusion policies, conceptions of ASD and pedagogical practices — and an integrative synthesis. Results indicate that normative advances were not matched, to the same extent, by effective school conditions. Insufficient initial teacher education, scarce resources and specialized support, curricular rigidity and attitudinal barriers persist. As strategies, the literature highlights curriculum adaptation, structured teaching, alternative communication, assistive technologies, play mediation, collaborative work and family participation. It is concluded that inclusion is not achieved by enrollment alone, but by structural, pedagogical and attitudinal reorganization of the school. The study is limited by its narrative design, without a systematic search, and by not assessing the methodological quality of the works. Further empirical research on classroom practices, teacher education and trajectories of students with ASD is recommended.

\vspace{1em}
\noindent\textbf{Keywords:} Inclusive education. Autism Spectrum Disorder. Early years. Teacher education. Pedagogical practices.
\endgroup

\clearpage

\begin{center}
\textbf{LISTA DE ABREVIATURAS E SIGLAS}
\end{center}

\begingroup
\setstretch{1.0}
\begin{description}
\item[ABNT] Associação Brasileira de Normas Técnicas
\item[CBEE] Congresso Brasileiro de Educação Especial
\item[CONEDU] Congresso Nacional de Educação
\item[DOI] Digital Object Identifier
\item[LBI] Lei Brasileira de Inclusão
\item[LDB] Lei de Diretrizes e Bases da Educação Nacional
\item[MEC] Ministério da Educação
\item[PNEEPEI] Política Nacional de Educação Especial na Perspectiva da Educação Inclusiva
\item[REIN] Revista Educação Inclusiva
\item[TEA] Transtorno do Espectro Autista
\item[UEPB] Universidade Estadual da Paraíba
\item[UFSM] Universidade Federal de Santa Maria
\end{description}
\endgroup

\clearpage
